\documentclass[12pt]{article}

\usepackage[utf8]{inputenc}
\usepackage[T1]{fontenc}
\usepackage[spanish, english]{babel}
\usepackage{amsmath}
\usepackage{booktabs}
\usepackage{graphicx}
\usepackage{geometry}
\usepackage{hyperref}
\usepackage{natbib}
\usepackage{longtable}

\geometry{margin=1in}
\graphicspath{{Paper/figures/}{figures/}}
\newcommand{\tablepath}{Paper/tables/}
\newcommand{\figurelang}{en}
\newif\ifpolicyreport
\policyreportfalse

\title{Firm Size, Wages, and Allocative Frictions in Colombia}
\author{Oliver Pardo and Jorge Orozco}
\date{\today}

\begin{document}

\maketitle

\begin{abstract}
Firm-size wage premia are well documented in advanced economies, but their
magnitude and interpretation in middle-income labor markets with large informal
sectors remain less settled. This paper documents the firm-size wage premium in
Colombia using harmonized microdata from the Gran Encuesta Integrada de Hogares
covering 2008 to 2025, excluding 2020. Employment is concentrated in solo work
and micro-firms, while real hourly wages rise sharply with firm size. After
controlling for gender, education, formality, and sector-year fixed effects,
workers in firms with 101 or more employees earn about 35 percent more than
otherwise comparable solo workers. The premium is positive throughout the sample
period, with no evidence of secular decline. Because the premium survives
formality controls, it is not simply a relabeling of the formal-informal wage
gap. Combined with the thin middle of the Colombian firm-size distribution, the
evidence points to a labor market in which firm scale is closely linked to wage
inequality, segmentation, and possible allocative frictions.
\end{abstract}


\section{Introduction}

Large firms tend to pay higher wages than small firms. This empirical regularity, usually known as the firm-size wage premium, is one of the classic facts in labor economics, and it often remains after conditioning on observed worker characteristics \citep{BrownMedoff1989,OiIdson1999,troske1999evidence}. Nonetheless, a firm-size wage premium conditional on observables should be interpreted carefully. It means that workers in larger firms earn more than observably similar workers in smaller firms, but it does not prove that moving a given worker to a larger firm would raise aggregate output, since unobserved worker quality or other omitted factors may still matter. At the same time, the persistence of the premium after controlling for observables is informative: it is consistent with allocative frictions even if it is not direct proof of it. Evidence from developing economies and matched employer--employee studies supports this cautious interpretation: firm-size effects often remain after stronger controls for worker heterogeneity, but their interpretation depends on the underlying source of the employer-associated wage gap \citep{SoderbomTealWambugu2005,reed2019large,AbowdKramarzMargolis1999,CardHeiningKline2013,Song2019,AlvarezBenguriaEngbomMoser2018,GerardLagosSeverniniCard2021,EngbomMoser2022}. The firm-size wage premium is therefore a diagnostic statistic rather than a causal estimate.

Our motivation to measure the firm-size wage premium in Colombia is both empirical and policy-oriented. A recent policy diagnosis argues that low-income households are disproportionately linked to own-account work, very small businesses, informality, and low-earning jobs, whereas higher-income households are more often connected to larger and more formal employers \citep{eslava2026crecimiento}. This diagnosis is suggestive, but the observed association between firm size and household income may reflect worker sorting or formality rather than firm scale itself. Nonetheless, the wage gradient associated with firm size has not been systematically documented in Colombia. Existing work has studied wage differentials by sector and firm characteristics in the formal labor market \citep{iregui2012diferenciales}, but less is known about the relationship between firm size and wages once solo work,
micro-firms, and informal employment are included. The central empirical challenge is therefore to document the firm-size wage gradient while distinguishing it from worker characteristics and the formal-informal wage gap.

This paper documents the firm-size wage premium in Colombia using harmonized microdata from the Gran Encuesta Integrada de Hogares (GEIH) for 2008--2019 and 2021--2025. The central questions are descriptive and econometric: do workers in larger firms earn higher real hourly labor income than comparable workers in solo work or smaller firms? How much of the raw gradient remains after accounting for observable worker characteristics, formality, and sector-year conditions? Has the premium widened or narrowed over time? Does the relationship between firm size and wages differ between formal and informal workers? The GEIH does not identify employers or observe marginal products, so the resulting estimates are conditional wage gaps, not causal effects of firm size. They are nevertheless informative because persistent conditional gaps in a labor market with concentrated micro-firm employment are a useful diagnostic of segmentation and possible allocative frictions.

We document five main facts. First, employment is concentrated at the bottom of the firm-size distribution: solo workers and firms with up to ten workers account for about two thirds of employment. Second, the raw wage gradient by firm size is large and widespread across sectors, education groups, gender, and formality status. Third, the gradient attenuates substantially after adding worker controls, sector-year fixed effects, and especially formality, but it does not disappear: in the full specification, workers in firms with 101 or more employees earn about 43 percent more than otherwise comparable solo workers. Fourth, the conditional premium is positive throughout the period, with no statistically conclusive evidence of either a secular decline or a systematic widening. Fifth, the formality status affects the firm-size wage premium. Formal workers earn more than informal workers within every firm-size category, but the conditional firm-size wage premium is much steeper among informal workers. Within formal employment, the broad size gradient is weak relative to formal solo workers, although workers in 101+ firms earn more than workers in intermediate-size formal firms. If the wage gaps are interpreted as a markers of allocative frictions rather than as pure sorting, the evidence points to three possible margins for efficiency gains: movement from informal to formal employment, scaling of informal micro-units, and, more selectively, movement or growth within the formal sector toward larger employers.

The rest of the paper is organized as follows. Section~2 describes the GEIH data, the harmonization of firm-size categories, and the construction of real hourly labor income. Section~3 presents the descriptive evidence on employment shares and wage gaps by firm size, formality, education, gender, sector, and year. Section~4 introduces the econometric specifications. Section~5 reports the main regression results, the dynamic estimates, and the formality interactions. Section~6 discusses interpretation and limitations, Section~7 draws policy implications, and Section~8 concludes.

\section{Data}

\subsection{Source}

The empirical analysis uses microdata from the Gran Encuesta Integrada de
Hogares (GEIH), Colombia's nationally representative labor force survey produced
by DANE. The current harmonized file covers the years 2008--2019 and 2021--2025.
The working sample contains 5,115,921 worker-year observations with valid real
hourly labor income and positive expansion weights. The baseline regression
sample contains 5,034,706 observations after requiring non-missing values for the
controls used in the econometric specification.

\subsection{Sample Construction}

The sample is restricted to employed workers with positive labor income, valid
weekly hours, and positive survey expansion factors. Monthly labor income is
taken from the GEIH labor income variable (\texttt{INGLABO}). We convert it to
annual labor income and divide by annual hours, computed as usual weekly hours
times 52:

\[
  w_{it}^{nominal} =
  \frac{12 \times INGLABO_{it}}{52 \times hours_{it}}.
\]

We then express hourly labor income in constant 2025 pesos using the December CPI
series encoded in the data construction scripts:

\[
  w_{it}^{real} = w_{it}^{nominal} \times
  \frac{CPI_{2025}}{CPI_t}.
\]

The main outcome is the logarithm of real hourly labor income,
\(\ln(w_{it}^{real})\). Observations with non-positive income, invalid hours, or
non-positive weights are excluded.

\subsection{Firm Size and Controls}

Firm size is measured using the worker-reported firm-size variable in GEIH. The
construction harmonizes the older \texttt{P6870} categories and the newer
\texttt{P3069} categories into the following ordered groups: Solo, 2-3, 4-5,
6-10, 11-19, 20-30, 31-50, 51-100, and 101+. When the newer questionnaire
separates 101-200 and 201+ workers, both categories are collapsed into 101+ to
maintain comparability with earlier years.

The control variables are also harmonized across survey changes. Education is
grouped into comparable categories: no education, preschool, primary, lower
secondary, upper secondary, and higher education. Formality is defined using
pension contribution status: workers who contribute to pensions are classified
as formal, while workers who do not contribute are classified as informal.
Economic activity is harmonized across CIIU revisions into broad sectors, which
allows the regressions to absorb sector-year shocks and changes in the sectoral
composition of employment.

Two caveats are important. First, GEIH is a household survey, not linked
employer-employee administrative data. Firm size is therefore reported by
workers, and the data do not identify firms. Second, the analysis documents
conditional wage differences by firm size; it does not identify the causal effect
of moving a worker from a small firm to a large firm.

\section{Descriptive Statistics}

The descriptive section is organized around the empirical facts that motivate the
regression analysis. We first define the working sample, then document where
workers are located in the firm-size distribution, how wages vary across that
distribution, and how much of the raw pattern is related to formality and other
observable wage gaps.

%\begin{table}[htbp]
  \centering
  \caption{Sample construction and coverage}
  \label{tab:sample-construction}
  \small
  \begin{tabular}{lp{0.58\textwidth}}
    \toprule
    Statistic & Value \\
    \midrule
    Harmonized worker-year file & 3,921,364 \\
    Valid real hourly income and positive weights & 3,921,364 \\
    Regression/descriptive sample & 3,852,322 \\
    Years covered & 2008--2025 (17 survey years; 2020 absent) \\
    Firm-size categories & Solo, 2-3, 4-5, 6-10, 11-19, 20-30, 31-50, 51-100, 101+ \\
    \bottomrule
  \end{tabular}
\end{table}


%Table~\ref{tab:sample-construction} defines the working sample and period coverage. The sample contains just over 5.0 million worker-year observations with valid wages, firm-size information, worker controls, and expansion weights over 17 survey years. Detailed employment, wage, and worker-composition tables are reported in Appendix~\ref{app:additional-descriptives}; the main text keeps the focus on the exhibits that structure the interpretation.

Figure~\ref{fig:income-mean-comparison} establishes the first raw fact behind the paper: mean real hourly income rises sharply with firm size. The gradient is visible in both endpoint years. In 2008, solo workers earned 5,431 COL 2025 pesos per hour on average, compared with 13,026 COL 2025 pesos per hour among workers in firms with 101 or more employees, a ratio of 2.4 to 1. By 2025, the corresponding means were 6,534 and 16,859 COL 2025 pesos per hour, a ratio of 2.6 to 1. %The middle categories also sit well above solo work in 2025: workers in firms with 6--10 employees earned 9,663 COL 2025 pesos per hour, and those in firms with 51--100 employees earned 13,217 COL 2025 pesos per hour. 
The bubble sizes show that these are not comparisons among negligible cells: solo work accounted for 34.7 percent of employment in 2025, while the 101+ category accounted for 25.2 percent.

\begin{figure}[htbp]
  \centering
  \includegraphics[width=0.95\textwidth]{fig65_\figurelang.png}
  \caption{Mean real hourly income by firm size}
  \label{fig:income-mean-comparison}
\end{figure}

Figure~\ref{fig:employment-share-comparison} establishes the second raw fact: employment is heavily concentrated at the bottom of the firm-size distribution. In 2025, solo workers accounted for 34.7 percent of employment, firms with 2--3 workers for 13.8 percent, firms with 4--5 workers for 7.0 percent, and firms with 6--10 workers for 5.8 percent. Taken together, solo work and firms with up to ten workers represented 61.2 percent of employment. By contrast, the middle of the distribution was thin: firms with 11--100 workers jointly accounted for 13.6 percent. The largest category grew substantially, from 18.8 percent of employment in 2008 to 25.2 percent in 2025, but this increase did not eliminate the large employment mass in solo work and micro-firms, which still covered almost three out of every five workers in 2025.

\begin{figure}[htbp]
  \centering
  \includegraphics[width=0.9\textwidth]{fig64_\figurelang.png}
  \caption{Employment distribution by firm size: 2008 and 2025}
  \label{fig:employment-share-comparison}
\end{figure}

Figure~\ref{fig:wage-gaps-2025} puts the firm-size gradient in context by comparing it with other raw remuneration gaps in 2025. The gender gap is small in these weighted means: men earn 0.6 percent below the overall mean and women 0.8 percent above it. The formal-informal gap is much larger. Formal workers earn 15,241 COL 2025 pesos per hour, 48.0 percent above the overall mean, while
informal workers earn 6,271 COL 2025 pesos per hour, 39.1 percent below it; the formal-informal wage ratio is therefore about 2.4 to 1, close to the 2.6-to-1 ratio between 101+ firms and solo work in
Figure~\ref{fig:income-mean-comparison}. Education and sector gaps are also large: workers with higher education earn 71.6 percent above the overall mean, while workers with preschool or less earn 57.9 percent below it; across activities, agriculture is 41.4 percent below the mean, whereas information and communication is 112.8 percent above it. The raw firm-size gap is therefore far
larger than the gender gap, comparable to the formality gap, and embedded in education and sectoral sorting. Appendix Table~\ref{tab:descriptive-wage-gaps-2025} reports the corresponding means, employment shares, and gaps.

\begin{figure}[!htbp]
  \centering
  \includegraphics[width=0.95\textwidth]{fig73_\figurelang.png}
  \caption{Mean real hourly income gaps by sex, formality, education, and sector, 2025; gaps are relative to the overall 2025 mean. The sector panel omits extraterritorial organizations.}
  \label{fig:wage-gaps-2025}
\end{figure}

Appendix Table~\ref{tab:descriptive-wage-groups-mean-2008-2025} reports mean real hourly
wages in 2008 and 2025 for all workers, by sex, and by formality status. The
table gives the exact magnitudes behind Figure~\ref{fig:income-mean-comparison}.
Three patterns stand out. First, the raw firm-size gradient was already large in
2008 and remained large in 2025: mean hourly income in the 101+ category was
about 2.4 times the solo-worker mean in 2008 and about 2.6 times the
solo-worker mean in 2025. Second, the gender gap in the aggregate means is small
relative to the cross-firm-size gradient. Third, the formal-informal wage gap is
large in both years, making formality a central part of the interpretation.
Appendix Table~\ref{tab:descriptive-wage-growth-2008-2025} reports the
corresponding wage changes between 2008 and 2025.

\ifpolicyreport
\begin{table}[htbp]
  \centering
  \caption{Real hourly labor income by firm size, 2025}
  \label{tab:descriptive-income-size-2025}
  \small
  \begin{tabular}{lrrrrrr}
    \toprule
    Firm size & Mean & P25 & Median & P75 & S.D. & Raw premium \\
    \midrule
    Solo & 8,351 & 3,606 & 5,769 & 8,942 & 11,807 & 0.0\% \\
    2-3 & 8,830 & 4,615 & 6,522 & 8,654 & 13,608 & 5.7\% \\
    4-5 & 9,503 & 5,440 & 6,923 & 8,654 & 16,211 & 13.8\% \\
    6-10 & 10,609 & 6,393 & 7,466 & 9,532 & 17,310 & 27.0\% \\
    11-19 & 12,346 & 6,844 & 7,867 & 11,276 & 18,438 & 47.8\% \\
    20-30 & 11,997 & 7,139 & 8,145 & 11,538 & 13,414 & 43.6\% \\
    31-50 & 13,428 & 7,141 & 8,515 & 13,516 & 17,573 & 60.8\% \\
    51-100 & 14,481 & 7,212 & 8,897 & 14,685 & 18,799 & 73.4\% \\
    101+ & 18,502 & 7,615 & 11,037 & 20,769 & 21,986 & 121.5\% \\
    \bottomrule
  \end{tabular}
  \vspace{0.3em}
  \begin{minipage}{0.95\textwidth}
  \footnotesize
  Notes: Incomes use only 2025 observations and are expressed as real hourly labor income in constant 2025 pesos. All statistics are weighted by GEIH expansion weights. The raw premium is the percent difference in the weighted mean relative to solo workers in 2025.
  \end{minipage}
\end{table}


\begin{table}[htbp]
  \centering
  \caption{Worker characteristics by firm size, 2025}
  \label{tab:descriptive-worker-characteristics-2025}
  \small
  \begin{tabular}{lrrrr}
    \toprule
    Firm size & Women & Formal & Higher education & Mean log wage \\
    \midrule
    Solo & 46.2\% & 14.1\% & 25.2\% & 8.639 \\
    2-3 & 42.4\% & 20.1\% & 26.1\% & 8.769 \\
    4-5 & 42.6\% & 32.2\% & 28.3\% & 8.879 \\
    6-10 & 44.0\% & 54.4\% & 36.3\% & 9.009 \\
    11-19 & 44.2\% & 70.4\% & 46.1\% & 9.128 \\
    20-30 & 45.1\% & 84.2\% & 49.5\% & 9.168 \\
    31-50 & 44.2\% & 91.9\% & 56.3\% & 9.254 \\
    51-100 & 46.8\% & 94.8\% & 58.4\% & 9.311 \\
    101+ & 47.5\% & 98.3\% & 67.6\% & 9.511 \\
    \bottomrule
  \end{tabular}
  \vspace{0.3em}
  \begin{minipage}{0.95\textwidth}
  \footnotesize
  Notes: Statistics use only 2025 observations and are weighted by GEIH expansion weights. Higher education corresponds to workers classified as superior or university educated.
  \end{minipage}
\end{table}


For the policy report, Tables~\ref{tab:descriptive-income-size-2025} and
\ref{tab:descriptive-worker-characteristics-2025} summarize the latest observed
labor-market structure in 2025.
\fi

\ifpolicyreport
\else
Figure~\ref{fig:income-formality-comparison} separates the endpoint comparison
by formality status. The bubble size represents the percentage of formal or
informal workers in each firm-size category in the corresponding year. Formal
workers earn more than informal workers and are much more concentrated in large
firms, but the wage gradient by firm size is visible within both formality
groups. This is why formality is a central mechanism but not the only object of
interest.

\begin{figure}[!htbp]
  \centering
  \includegraphics[width=0.95\textwidth]{fig69_\figurelang.png}
  \caption{Mean real hourly income by firm size and formality status, 2008 and 2025; dashed hollow markers denote 2008, solid filled markers denote 2025, and bubble size shows each firm-size category's employment share within the corresponding formality-year}
  \label{fig:income-formality-comparison}
\end{figure}

Figure~\ref{fig:income-formality-comparison} also shows that the firm-size gradient is not simply a formal-informal contrast. Among informal workers, mean hourly income still rises from 5,677 pesos for solo workers to 12,813 pesos in firms with 101 or more workers in 2025. Among formal workers, however, the profile is closer to a U-shape: formal solo workers earned 15,990 pesos per hour in 2025, formal workers in firms with 4--5 workers earned 12,036 pesos, and formal workers in firms with 101 or more workers earned 16,944 pesos. A plausible interpretation is that formal solo work is a selected category that includes independent professionals, such as lawyers, physicians, and consultants. This interpretation is consistent with the high earnings of formal solo workers and with their relatively small employment share within formal employment.

Figures~\ref{fig:income-sex-comparison} and
\ref{fig:income-education-comparison} extend the endpoint comparison to sex and
education groups.

\begin{figure}[!htbp]
  \centering
  \includegraphics[width=0.95\textwidth]{fig70_\figurelang.png}
  \caption{Mean real hourly income by firm size and sex, 2008 and 2025; dashed hollow markers denote 2008, solid filled markers denote 2025, and bubble size shows each firm-size category's employment share within the corresponding sex-year}
  \label{fig:income-sex-comparison}
\end{figure}

Figure~\ref{fig:income-sex-comparison} shows that the raw firm-size gradient is
similar for men and women. In 2025, male solo workers earned 6,816 pesos per
hour, compared with 16,929 pesos in firms with 101 or more workers; for women,
the corresponding means were 6,133 and 16,777 pesos. The large-firm advantage is
therefore visible within both sexes. The figure also shows that the gender gap
within a given firm-size category is small relative to the gap between solo work
and large-firm employment.

\begin{figure}[!htbp]
  \centering
  \includegraphics[width=\textwidth]{fig72_\figurelang.png}
  \caption{Mean real hourly income by firm size and education group, 2008 and 2025; dashed hollow markers denote 2008, solid filled markers denote 2025, and bubble size shows each firm-size category's employment share within the corresponding education-year}
  \label{fig:income-education-comparison}
\end{figure}

Figure~\ref{fig:income-education-comparison} is especially informative because
it separates two mechanisms that are otherwise mixed together. First, education
composition matters. In 2025, workers with preschool or less, primary, and
secondary education were much more concentrated in solo work than workers with
higher education. Solo work accounted for 65.2 percent of employment among
workers with preschool or less and 55.5 percent among workers with primary
education, compared with only 17.3 percent among workers with higher education.
Conversely, the 101+ category accounted for 47.3 percent of workers with higher
education but only 1.4 percent of workers with preschool or less and 4.9 percent
of workers with primary education. This sorting helps explain why the raw
firm-size wage gradient is so steep.

Second, education composition does not exhaust the pattern. Within every
education group, workers in the largest firms earned more than solo workers in
2025. The ratio between mean hourly income in 101+ firms and solo work was 2.1
among workers with preschool or less, 1.8 among workers with primary education,
1.5 among workers with secondary education, and 1.7 among workers with higher
education. At the same time, the higher-education panel has a distinctive shape:
solo workers with higher education earned more than workers in firms with 2--3
employees, before the gradient rose toward large firms. This is consistent with
the interpretation that some solo work among highly educated workers reflects
professional self-employment rather than only low-productivity informality. The
analogous comparison by sector is reported in
Appendix~\ref{app:additional-descriptives}.
\fi

\ifpolicyreport
\noindent Figures~\ref{fig:policy-employment-income-2025} and
\ref{fig:policy-composition-2025} summarize the 2025 policy snapshot, combining
employment, wage, formality, and education gradients by firm size.

\begin{figure}[htbp]
  \centering
  \includegraphics[width=0.9\textwidth]{fig66_\figurelang.png}
  \caption{Employment and mean hourly income by firm size in 2025}
  \label{fig:policy-employment-income-2025}
\end{figure}

\begin{figure}[htbp]
  \centering
  \includegraphics[width=0.9\textwidth]{fig67_\figurelang.png}
  \caption{Formality and higher education by firm size in 2025}
  \label{fig:policy-composition-2025}
\end{figure}
\fi

Taken together, the descriptive evidence motivates the regression analysis in
four ways. First, employment is concentrated in low-scale units: in the pooled
regression sample, solo workers and firms with up to ten workers account for
65.3 percent of weighted employment; in 2025 they still account for 61.3
percent. Second, the middle of the firm-size distribution is thin: firms with
11--100 workers account for only 12.3 percent of pooled employment. Third, the
largest harmonized category is economically important: firms with 101 or more
workers account for 22.5 percent of pooled employment. Fourth, the raw wage
gradient is steep, but larger firms also employ more formal and more educated
workers. The econometric exercise below therefore asks how much of the raw
gradient survives after controlling for observable worker composition,
formality, and sector-year conditions.

%\clearpage

\section{Methodology}

The main econometric exercise estimates the conditional firm-size wage premium
in a log-wage model. The baseline specification is:

\[
  \ln(w_{ist}) =
  \alpha
  + \sum_{g \neq g_0} \beta_g \mathbf{1}\{Size_{ist}=g\}
  + X_{ist}'\gamma
  + \lambda_{s \times t}
  + \varepsilon_{ist},
\]

where \(w_{ist}\) is real hourly labor income for worker \(i\) in sector \(s\) in year \(t\), \(Size_{ist}\) is the harmonized firm-size category, and \(g_0\) is the reference group. We use solo workers as the reference group. This choice makes the coefficient path directly interpretable as the wage premium associated with moving from individual work to increasingly larger firms. Because solo workers may combine self-employment and very small-scale wage employment, we interpret the estimates as conditional gaps relative to the bottom of the firm-size distribution rather than as pure firm effects.

The vector \(X_{ist}\) includes a female-worker indicator, age, age squared,
education dummies, and a formal-employment indicator. The fixed effect
\(\lambda_{s \times t}\)
absorbs sector-year conditions, allowing each broad sector to have its own wage
level in each year. This is important because large firms are not randomly
distributed across sectors, and sectoral shocks may be correlated with both wages
and firm size. The model is estimated with GEIH expansion weights, and standard
errors are clustered by sector.

City fixed effects would be a useful additional robustness check because they
would absorb persistent differences across local labor markets. The harmonized
file used for the regressions does not yet preserve a consistent city or
municipality identifier, but it does retain department/Bogot\'a. Appendix
Table~\ref{tab:firm-size-local-geography-robustness} therefore reports a
stringent local robustness exercise that absorbs sector-department-year fixed
effects.

In the results table, we report four nested specifications. The first column
includes only firm-size indicators. The second adds sector-year fixed effects.
The third adds gender, age, age squared, and education controls, but excludes
formality. The fourth adds formality. Comparing the third and fourth columns
shows how much of the
firm-size gradient is absorbed by differences in formal employment across firm
sizes.

The coefficients \(\beta_g\) are log-point differences relative to solo workers.
For interpretation, we transform them into percentage premiums:

\[
  Premium_g = 100 \times \left(\exp(\hat{\beta}_g)-1\right).
\]

\subsection{Firm-Size Nonlinearities by Formality Status}

The descriptive evidence suggests that the relationship between firm size and
hourly income may differ between formal and informal workers. To allow for this
nonlinearity, we also estimate a specification that interacts all firm-size
categories with formal employment:

\[
  \ln(w_{ist}) =
  \alpha
  + \rho Formal_{ist}
  + \sum_{g \neq g_0} \beta_g \mathbf{1}\{Size_{ist}=g\}
  + \sum_{g \neq g_0} \theta_g
    \mathbf{1}\{Size_{ist}=g\} Formal_{ist}
  + Z_{ist}'\gamma
  + \lambda_{s \times t}
  + \varepsilon_{ist}.
\]

Here \(Z_{ist}\) includes the worker controls other than formality. The
coefficients \(\beta_g\) measure the firm-size premium among informal workers,
relative to informal solo workers. The corresponding premium among formal
workers is \(\beta_g+\theta_g\), relative to formal solo workers. This
specification is deliberately nonparametric in firm size: it does not impose a
linear, quadratic, or monotonic relationship between firm size and wages.

\subsection{Dynamic Premiums and Trend Test}

The same framework can be extended to study the evolution of the premium over
time by interacting firm-size categories with year indicators:

\[
  \ln(w_{ist}) =
  \alpha
  + \sum_{\tau}\sum_{g \neq g_0}
    \beta_{g\tau}\mathbf{1}\{Size_{ist}=g\}\mathbf{1}\{Year_{t}=\tau\}
  + X_{ist}'\gamma
  + \lambda_{s \times t}
  + \varepsilon_{ist}.
\]

This dynamic specification produces a sequence of firm-size premiums by year and
can be used to assess whether the premium has widened, narrowed, or remained
stable. We also estimate a more parsimonious trend specification:

\[
  \ln(w_{ist}) =
  \alpha
  + \sum_{g \neq g_0} \beta_g\mathbf{1}\{Size_{ist}=g\}
  + \sum_{g \neq g_0} \delta_g
    \mathbf{1}\{Size_{ist}=g\}Trend_t
  + X_{ist}'\gamma
  + \lambda_{s \times t}
  + \varepsilon_{ist},
\]

where \(Trend_t\) is a linear year trend normalized to zero in 2008. The
parameter \(\delta_g\) tests whether the wage premium for firm-size group \(g\)
increased over time relative to solo workers. We report one-sided tests of
\(H_0:\delta_g=0\) against \(H_1:\delta_g>0\), because the empirical question is
whether the premium widened over the period.

\section{Results}

\subsection{Main Results}

Table~\ref{tab:firm-size-wage-premium-regressions} reports the regression results from the main specification in Equation (\ref{eq:main}). Column (1) shows the raw log-wage gradient by firm size. The coefficient for firms with 101 or more workers is 1.019 log points, consistent with the large raw wage premium shown in the descriptive statistics. Column (2) adds sector-year fixed effects, reducing the coefficient for 101+ firms to 0.819 log points. Column (3) adds gender, age, age squared, and education controls, lowering the coefficient to 0.598. Column (4) adds formality, reducing the coefficient further to 0.359. Thus, much of the raw firm-size wage gap is associated with sector-year conditions and observable differences in worker and job characteristics, including gender, age, education, and formality. Nonetheless, a positive and statistically significant conditional premium remains after controlling for these observables, suggesting that barriers to firm growth may matter for workers' wages. If larger firms are more productive or better able to share rents with workers, then obstacles that prevent small firms from scaling up, or workers from moving into larger firms, can depress earnings and contribute to wage inequality and labor-market inefficiencies.

\begin{table}[htbp]
  \centering
  \caption{Firm-size wage premium regressions: the role of formality}
  \label{tab:firm-size-wage-premium-regressions}
  \small
  \begin{tabular}{lcccc}
    \toprule
    & \multicolumn{4}{c}{Dependent variable: log real hourly labor income} \\
    \cmidrule(lr){2-5}
    & (1) & (2) & (3) & (4) \\
    \midrule
    Firm size: 2-3 & 0.139*** & 0.193*** & 0.168*** & 0.145*** \\
     & (0.033) & (0.029) & (0.027) & (0.029) \\
    Firm size: 4-5 & 0.287*** & 0.325*** & 0.293*** & 0.234*** \\
     & (0.036) & (0.034) & (0.035) & (0.044) \\
    Firm size: 6-10 & 0.390*** & 0.402*** & 0.337*** & 0.234*** \\
     & (0.034) & (0.038) & (0.041) & (0.052) \\
    Firm size: 11-19 & 0.492*** & 0.468*** & 0.363*** & 0.224*** \\
     & (0.039) & (0.043) & (0.044) & (0.050) \\
    Firm size: 20-30 & 0.529*** & 0.479*** & 0.356*** & 0.190*** \\
     & (0.042) & (0.047) & (0.046) & (0.050) \\
    Firm size: 31-50 & 0.592*** & 0.521*** & 0.378*** & 0.195*** \\
     & (0.052) & (0.051) & (0.048) & (0.049) \\
    Firm size: 51-100 & 0.657*** & 0.567*** & 0.407*** & 0.215*** \\
     & (0.061) & (0.053) & (0.047) & (0.047) \\
    Firm size: 101+ & 0.843*** & 0.672*** & 0.480*** & 0.282*** \\
     & (0.094) & (0.069) & (0.063) & (0.064) \\
    \midrule
    Gender, age, and education controls & No & No & Yes & Yes \\
    Formality control & No & No & No & Yes \\
    Sector-year fixed effects & No & Yes & Yes & Yes \\
    Observations & 3,852,322 & 3,852,322 & 3,852,322 & 3,852,322 \\
    $R^2$ & 0.162 & 0.221 & 0.353 & 0.362 \\
    \bottomrule
  \end{tabular}
  \vspace{0.3em}
  \begin{minipage}{0.95\textwidth}
  \footnotesize
  Notes: The omitted category is solo workers. All columns use the same estimation sample and GEIH expansion weights. Standard errors, clustered by sector, are reported in parentheses. Worker controls include a female-worker indicator, age, age squared, and education dummies. Column (4) adds labor formality. Significance levels: * $p<0.10$, ** $p<0.05$, *** $p<0.01$.
  \end{minipage}
\end{table}


%The comparison across columns is central to the interpretation. The raw premium documents the inequality that workers actually face across the firm-size distribution. The reduction after adding education, age, gender, sector-year fixed effects, and formality shows that part of the gap reflects sorting and segmentation: larger firms employ a more formal and more educated workforce and operate in different sector-year environments. The remaining conditional premium does not prove that larger firms causally raise wages, but it is consistent with firm-level wage-setting differences, rent sharing, internal labor markets, technology, productivity, or unobserved worker-firm matching. In a labor market where most workers are concentrated in solo work and micro-firms, this remaining premium is economically meaningful even after the controls absorb a large share of the raw gap.

%Figure~\ref{fig:firm-size-premium-attenuation} summarizes this attenuation for selected firm-size categories. The figure shows that the raw gap is large, but that sector-year conditions, worker composition, and especially formality absorb an important share of the difference between small and large firms.

%\begin{figure}[htbp]
%  \centering
%  \includegraphics[width=0.9\textwidth]{fig68_\figurelang.png}
%  \caption{Attenuation of the firm-size wage premium across specifications}
%  \label{fig:firm-size-premium-attenuation}
%\end{figure}

Figure~\ref{fig:firm-size-wage-premium} transforms the coefficients from Column (4) in Table (\ref{tab:firm-size-wage-premium-regressions}) into percentage premiums. Relative to solo workers, the conditional premium is 20.7 percent for firms with 2--3 workers, 36.0 percent for firms with 4--5 workers, 37.3 percent for firms with 6--10 workers, roughly 31--35
percent for intermediate size categories, and 43.2 percent for firms with 101 or more workers. All confidence intervals exclude zero.

\begin{figure}[htbp]
  \centering
  \includegraphics[width=0.9\textwidth]{fig57_\figurelang.png}
  \caption{Conditional firm-size wage premium relative to solo workers}
  \label{fig:firm-size-wage-premium}
\end{figure}

\subsection{Dynamic results}

Figure~\ref{fig:firm-size-wage-premium-over-time} estimates the full specification separately by interacting firm-size categories with year indicators. The conditional premium is positive throughout the sample period in all firm-size categories. The largest firms display the highest premiums in most years, although the confidence intervals are wider for these categories.

\begin{figure}[htbp]
  \centering
  \includegraphics[width=0.95\textwidth]{fig59_\figurelang.png}
  \caption{Conditional firm-size wage premium over time}
  \label{fig:firm-size-wage-premium-over-time}
\end{figure}

Table~\ref{tab:firm-size-premium-trend-test} formally tests whether the premium changed over time using a linear trend interacted with firm-size categories. The estimated trend coefficients are positive for all size categories. The implied increase between 2008 and 2025 ranges from 3.2 percent for firms with 20--30 workers to 7.9 percent for firms with 101 or more workers. However, the two-sided tests do not reject a zero trend at the 5 percent level. Thus, the current evidence supports a positive conditional firm-size wage premium, but the evidence that this premium has changed over time is suggestive rather than statistically conclusive. Since all estimated trend coefficients are positive, the suggestive pattern points toward widening rather than narrowing, but the bilateral tests do not support a statistically conclusive trend.

The dynamic exercise also clarifies what it would mean for the firm-size wage gradient to steepen or flatten over time. A steepening premium would imply that access to large-firm employment has become increasingly important for wages, raising concerns about inequality and the allocation of workers across the productive structure. A flattening premium would be ambiguous: it could reflect improving wages and productivity among smaller firms, but it could also reflect weaker wage growth among larger firms. For this reason, the policy-relevant
object is not only the gap between small and large firms, but the joint evolution of wage levels, employment shares, formality, and productivity across firm-size categories.

\begin{table}[htbp]
  \centering
  \caption{Testing whether the firm-size wage premium changed over time}
  \label{tab:firm-size-premium-trend-test}
  \small
  \begin{tabular}{lcccc}
    \toprule
    Firm size & Annual trend & S.E. & $p$-value & Implied 2008--2025 change \\
    \midrule
    2-3 & 0.0024 & (0.0029) & 0.409 & 4.3\% \\
    4-5 & 0.0005 & (0.0039) & 0.909 & 0.8\% \\
    6-10 & 0.0014 & (0.0036) & 0.693 & 2.5\% \\
    11-19 & 0.0022 & (0.0043) & 0.610 & 3.9\% \\
    20-30 & 0.0023 & (0.0042) & 0.595 & 3.9\% \\
    31-50 & 0.0029 & (0.0046) & 0.533 & 5.1\% \\
    51-100 & 0.0037 & (0.0042) & 0.393 & 6.5\% \\
    101+ & 0.0054 & (0.0047) & 0.268 & 9.7\% \\
    \bottomrule
  \end{tabular}
  \vspace{0.3em}
  \begin{minipage}{0.95\textwidth}
  \footnotesize
  Notes: The annual trend is the coefficient on the interaction between firm-size category and a linear year trend, with solo workers as the omitted category. The specification controls for gender, age, age squared, education, formality, and sector-year fixed effects, and uses GEIH expansion weights. Standard errors are clustered by sector. The $p$-value corresponds to the two-sided test that the trend differs from zero. The final column reports $100\times[\exp(17\hat{\delta})-1]$, the implied change in the firm-size wage ratio between 2008 and 2025.
  \end{minipage}
\end{table}


\subsection{Formality non-linear results}

Figure~\ref{fig:firm-size-formality-premium} asks whether the nonlinear pattern seen in Figure~\ref{fig:income-formality-comparison} survives after controlling for worker composition and sector-year conditions. The figure reports firm-size premiums relative to solo workers within each formality group. Among informal workers, the conditional premium is positive and statistically significant in every firm-size category: it is 21.7 percent in firms with 2--3 workers, rises to about 48--51 percent across most categories from 6 to 100 workers, and remains 40.0 percent in firms with 101 or more workers. Among formal workers, the within-formality gradient is much flatter and less precisely estimated: the estimated premium ranges from 2.1 to 7.0 percent for most categories and is 16.5 percent for firms with 101 or more workers, but the confidence intervals include zero throughout when each category is compared with formal solo workers. This does not mean, however, that there is no wage structure within formal employment. The U-shape in Figure~\ref{fig:income-formality-comparison} points to a different comparison: formal workers in 101+ firms versus formal workers in intermediate-size firms. Direct contrasts from the interacted model show that formal workers in 101+ firms earn between 8.9 and 14.2 percent more than formal workers in every smaller non-solo category, and all of these contrasts are statistically significant. Table~\ref{tab:formal-101-contrasts} reports these contrasts explicitly. Thus, the estimates are consistent with a heterogeneous and selected formal solo group coexisting with a large-firm premium relative to
the middle of the formal firm-size distribution. The informal sector, by contrast, displays a robust conditional firm-size gradient relative to solo work.

\begin{figure}[htbp]
  \centering
  \includegraphics[width=0.95\textwidth]{fig74_\figurelang.png}
  \caption{Adjusted firm-size wage premium by formality status}
  \label{fig:firm-size-formality-premium}
\end{figure}

Figure~\ref{fig:adjusted-formality-endpoints} provides the closest
regression-adjusted counterpart to
Figure~\ref{fig:income-formality-comparison}. The y-axis is the adjusted percent
premium relative to solo workers within the same formality-year group, rather
than an adjusted wage level. This normalization avoids choosing arbitrary
covariate values and makes the comparison with the solo-worker benchmark
explicit. Among formal workers, the raw U-shape largely disappears after
controls: in 2025, the adjusted estimates for all categories below 101+ range
from -3.7 to 0.6 percent, and the 101+ estimate is 10.5 percent with a
confidence interval that includes zero. The 2008 formal profile is similar.
Among informal workers, by contrast, the adjusted gradient remains steep and
statistically significant: in 2025, the premium rises from 24.3 percent in firms
with 2--3 workers to 53.6 percent in firms with 101 or more workers. Thus, the
descriptive U-shape among formal workers appears to be driven mainly by worker
composition within formal solo work, while the
informal firm-size premium survives the controls. This conclusion concerns the
comparison with solo workers; Table~\ref{tab:formal-101-contrasts} turns to the
separate comparison between large formal firms and the middle of the formal
firm-size distribution.

\begin{figure}[htbp]
  \centering
  \includegraphics[width=0.95\textwidth]{fig76_\figurelang.png}
  \caption{Adjusted firm-size wage premium by formality status, 2008 and 2025}
  \label{fig:adjusted-formality-endpoints}
\end{figure}

\begin{table}[htbp]
  \centering
  \caption{Formal-worker large-firm wage premium relative to other formal firm-size categories}
  \label{tab:formal-101-contrasts}
  \small
  \begin{tabular}{lccc}
    \toprule
    Reference category & Premium & 95\% CI & $p$-value \\
    \midrule
    Solo workers & 12.4\% & [-6.5, 35.1] & 0.213 \\
    2-3 & 12.5\% & [4.5, 21.0] & 0.002 \\
    4-5 & 12.9\% & [6.3, 19.9] & $<0.001$ \\
    6-10 & 13.1\% & [6.1, 20.5] & $<0.001$ \\
    11-19 & 11.3\% & [2.7, 20.6] & 0.009 \\
    20-30 & 13.3\% & [4.4, 22.9] & 0.003 \\
    31-50 & 11.2\% & [3.3, 19.7] & 0.005 \\
    51-100 & 8.1\% & [1.6, 15.0] & 0.013 \\
    \bottomrule
  \end{tabular}
  \vspace{0.3em}
  \begin{minipage}{0.92\textwidth}
  \footnotesize
  Notes: Each row compares formal workers in firms with 101 or more workers with formal workers in the reference firm-size category. Estimates come from the firm-size-by-formality specification with gender, age, age squared, education, and sector-year fixed effects. Premiums are computed as $100\times[\exp(\hat{\beta})-1]$ from the relevant linear contrast.
  \end{minipage}
\end{table}


Figure~\ref{fig:large-firm-formality-premium-over-time} examines whether the
pooled formal-worker result masks time variation. The figure plots the adjusted
premium for firms with 101 or more workers relative to solo workers, separately
within formal and informal employment in each year. The evidence does not show a
reappearance of a statistically significant large-firm premium among formal
workers in the most recent years. The formal 101+ premium is positive and
significant in several earlier years, especially 2009--2013 and 2017, but after
2021 the point estimates are smaller and the confidence intervals include zero;
in 2025, the estimate is 10.5 percent with a confidence interval from -8.2 to
33.0 percent. By contrast, the informal 101+ premium is positive and significant
throughout the period, rising from 39.4 percent in 2008 to 53.6 percent in
2025.

\begin{figure}[htbp]
  \centering
  \includegraphics[width=0.95\textwidth]{fig75_\figurelang.png}
  \caption{Adjusted large-firm wage premium by formality status over time; 2020 is omitted}
  \label{fig:large-firm-formality-premium-over-time}
\end{figure}



\subsection{Interpretation: Inequality, Efficiency, and Misallocation}

The results connect the firm-size wage premium to three related issues. The
first is wage inequality. Since workers in larger firms earn more, the allocation
of workers across firm sizes affects the overall wage distribution. The second is
labor-market segmentation. The strong role of formality in reducing the
coefficient for large firms suggests that the firm-size gradient is partly a
formal-informal gradient. The third is efficiency. If larger firms are more
productive and the wage premium partly reflects higher marginal products or
rents shared with workers, then the large employment mass in solo work and
micro-firms may indicate that workers and firms are not moving freely toward
higher-productivity matches.

This last interpretation should be treated as a hypothesis rather than as a
direct conclusion. The GEIH does not observe firms, capital, value added, or
marginal products, so the estimates cannot establish misallocation in the strict
sense of \citet{hsieh2009misallocation}. What the paper documents is a set of
facts consistent with a misallocation concern: employment is heavily concentrated
in very small units, larger firms pay substantially more, and a positive premium
remains after conditioning on observed worker characteristics, formality, and
sector-year fixed effects. Establishing the productivity channel directly would
require matched employer-employee data or firm-level information on output,
capital, and employment.

\subsection{Limitations}

Several limitations qualify the interpretation. First, firm size is reported by
workers in a household survey, and the data do not identify employers. Second,
the analysis cannot separate worker sorting from firm wage-setting: a remaining
conditional premium may reflect unobserved worker ability, firm rents, internal
labor markets, productivity differences, or match quality. Third, formality is
both a control and a potential mechanism. Adding it to the regression is useful
because it shows how much of the raw gradient is associated with the formal
segment of the labor market, but it may also absorb part of the channel through
which larger firms pay more. Finally, because the GEIH does not observe value
added, capital, or marginal products, the evidence should be read as suggestive
of allocative frictions rather than as a direct test of misallocation.

\section{Discussion}

Our central finding is a sizable firm-size wage premium that survives controls
for gender, age, age squared, education, formality, and sector-year conditions. While the
magnitude of the premium is broadly consistent with the existing Latin American
evidence on employer-side wage differences \citep{Schaffner1998,
AlvarezBenguriaEngbomMoser2018,GerardLagosSeverniniCard2021}, the persistence
of the premium after controlling for formality is the more substantively
informative feature of the results. This section discusses what such a premium
can and cannot tell us about the underlying mechanisms.

A long-standing concern in the developing-country literature is that the
apparent firm-size premium is, at least in part, a relabeling of the
formal-informal wage gap. In dual labor-market models in the tradition of
\citet{HarrisTodaro1970} and \citet{Fields1975}, large firms are
disproportionately formal, and formality itself carries a wage premium through
some combination of higher productivity, mandated benefits, payroll taxation,
and selection. From this perspective, much of what is conventionally labeled a
``firm-size'' premium may simply reflect the boundary between the formal and
informal sectors rather than any return to firm scale per se, since formal firms
also tend to be larger.

By showing that the premium remains after conditioning on formality, our
estimates indicate that this dual-market channel is not the whole story. Even
among workers with the same formality status, larger employers pay more. In this
respect, our findings reinforce \citet{Schaffner1998} and complement the
matched employer-employee evidence from Brazil
\citep{AlvarezBenguriaEngbomMoser2018,GerardLagosSeverniniCard2021,
EngbomMoser2022}, which points toward employer-side pay determination as a
first-order feature of Latin American labor markets. Related firm-to-firm
evidence from Costa Rica shows that productive linkages with multinationals can
shift domestic firm performance and employment
\citep{AlfaroUrenaManeliciVasquez2022}.

What our controls cannot do is distinguish among other leading explanations for
the premium. As \citet{BrownMedoff1989} and \citet{OiIdson1999} emphasized for
the United States, a premium that survives observable controls is consistent
with several non-exclusive mechanisms. The first is unobserved worker quality:
gender, schooling, and sector capture only part of productive heterogeneity,
leaving non-cognitive skills, school quality, and unmeasured ability bundled
into the firm-size coefficient. At the same time, the available developing-
country evidence cautions against treating sorting as the whole explanation:
\citet{SoderbomTealWambugu2005} show that the firm-size effect remains
substantial in Kenya and Ghana even after controlling for individual fixed
effects.

A second possibility is rent sharing. Larger firms tend to be more productive,
more capital-intensive, and more exposed to product-market power, and they may
share part of those rents with workers through bargaining, fair-wage norms, or
hold-up considerations. The matched-data literature beginning with
\citet{AbowdKramarzMargolis1999} and \citet{CardHeiningKline2013}, and extended
to the United States by \citet{Song2019}, has shown that firm-specific pay
premia account for a quantitatively important share of wage variation in
advanced economies; the Brazilian evidence cited above extends this conclusion
to Latin America. A third possibility is efficiency-wage or monitoring-cost
reasoning, in which larger firms face greater difficulty observing individual
effort and pay above-market wages to elicit it. Because our identification rests
on observables, we cannot adjudicate among these mechanisms, and we do not
interpret the surviving coefficient as a pure rent-sharing parameter.

%One reading we explicitly do not draw is a direct monopsony test. The recent literature on firm wage setting emphasizes that employers may have wage-setting power, but monopsony concerns markdowns relative to marginal products and firm-specific labor-supply elasticities, neither of which is observed in the GEIH \citep{Card2022}. A positive, robust firm-size premium is therefore compatible with monopsonistic labor markets, but it is not evidence for or against monopsony by itself. In this paper, the premium is more directly informative about productivity-driven pay differences, rent sharing, worker sorting, and segmentation, consistent with the firm-centric view of wage determination that has emerged in the recent literature.

Several limitations follow from this discussion. Because our identification rests on observables, we cannot separate unobserved worker quality from the effect of firm scale. Doing so would require either matched employer-employee data with worker mobility, in the spirit of \citet{AbowdKramarzMargolis1999} and \citet{CardHeiningKline2013}, or a source of plausibly exogenous variation in firm-size exposure. Similarly, while controlling for formality rules out a purely mechanical formal-informal explanation, it does not rule out firm characteristics correlated with both size and pay-setting capacity, such as capital intensity, exporter status, or foreign ownership. These extensions are natural next steps. In the meantime, our results contribute to a growing body of evidence that the firm-size wage premium in Latin America is significant, robust to standard observable controls, and not reducible to the formal-informal margin.

The evidence does not imply that policy should simply push workers into large firms or favor large firms as such. The paper does not observe firm productivity or worker-firm matches, so it cannot prove that the wage premium is a causal return to scale. The policy relevance of the results is more precise: a large share of Colombian employment remains in solo work and micro-firms, the middle of the firm-size distribution is thin, and larger firms pay higher wages even after conditioning on observed worker characteristics, formality, and sector-year conditions. This combination is a warning sign of labor-market segmentation and possible allocative frictions.

The policy implication is that policies should avoid preserving smallness as an objective in itself. Support for small firms is most valuable when it helps productive firms grow, formalize, adopt better technologies, and connect to larger markets. By contrast, policies that subsidize informality or create thresholds that make growth costly may reinforce the same employment structure documented in the paper: many workers in very small units, few workers in the middle, and a large wage gap between low-scale and large-scale employment.
\section{Policy Implications}

The evidence does not imply that policy should simply push workers into large
firms or favor large firms as such. The paper does not observe firm productivity
or worker-firm matches, so it cannot prove that the wage premium is a causal
return to scale. The policy relevance of the results is more precise: a large
share of Colombian employment remains in solo work and micro-firms, the middle
of the firm-size distribution is thin, and larger firms pay higher wages even
after conditioning on observed worker characteristics, formality, and
sector-year conditions. This combination is a warning sign of labor-market
segmentation and possible allocative frictions.

The first implication is that policies should avoid preserving smallness as an
objective in itself. Support for small firms is most valuable when it helps
productive firms grow, formalize, adopt better technologies, and connect to
larger markets. By contrast, policies that subsidize informality or create
thresholds that make growth costly may reinforce the same employment structure
documented in the paper: many workers in very small units, few workers in the
middle, and a large wage gap between low-scale and large-scale employment.

The second implication is that formalization should be treated as a mechanism,
not only as an outcome. The attenuation of the firm-size coefficient after
adding formality shows that the firm-size gradient is strongly connected to the
formal segment of the labor market. This supports policies that lower the fixed
costs of formal operation, simplify compliance, improve enforcement where evasion
is strategic, and make social protection less dependent on remaining outside
formal employment. At the same time, because the premium survives controlling
for formality, formalization alone is unlikely to eliminate the wage gradient.

The third implication concerns worker mobility. If large and growing firms offer
better wage opportunities, workers need credible paths into those jobs. This
points to training, job-search intermediation, certification, childcare and
transport policies, and regional mobility tools that reduce barriers between
workers in low-scale activities and firms with better matches. The descriptive
wage gaps by education and sector show why these policies matter: firm size is
not the only axis of inequality, and workers are sorted across schooling groups
and sectors before firm size enters the analysis.

Finally, the results should be linked to productivity policy. The premium is
consistent with, but does not prove, a misallocation story. A stronger test would
combine household-survey wage patterns with firm-level measures of value added,
capital, employment, and productivity. In the meantime, the evidence supports a
policy agenda focused on removing barriers to productive firm growth, improving
the quality of worker-firm matches, and helping workers move from low-scale
employment toward more productive, formal, and scalable firms. This is closely
aligned with the broader growth agenda for Colombia emphasized by
\citet{eslava2026crecimiento}: the objective is not firm size per se, but the
reorganization of production toward higher-productivity employers and better
jobs.

\section{Conclusion}

This paper documents a clear firm-size wage premium in Colombia. The descriptive evidence shows that real hourly wages rise sharply with firm size, while employment remains concentrated in solo work and micro-firms. Solo workers and firms with up to ten workers account for about two
thirds of employment, while the middle of the firm-size distribution remains thin.

The regression results show that a substantial conditional premium remains after controlling for gender, age, age squared, education, formality, and sector-year fixed effects. In the full specification, workers in firms with 101 or more employees earn about 43 percent more than otherwise comparable solo workers. At the same time, the large reduction in the coefficient as controls are added shows that worker composition, sector-year structure, and formality explain much but not all of the wage gradient.

The formal-informal results sharpen this interpretation. Formal workers earn more than informal workers within every firm-size category, so formality remains a large wage-level margin. However, the conditional wage gradient by firm size is
much steeper among informal workers than among formal workers. Within formal employment, the broad size gradient is weak when formal solo workers are the reference group, although workers in 101+ firms earn more than workers in intermediate-size formal firms. If the firm-size premium is interpreted as a
marker of allocative frictions rather than as pure unobserved sorting, the evidence therefore points to three possible margins for efficiency gains: moving workers and jobs from informality to formality, scaling informal micro-units into larger and more organized firms, and, more selectively, growth or movement within the formal sector from intermediate-size firms toward larger employers.

The central interpretation is therefore deliberately cautious. The firm-size wage premium is a marker of wage inequality and labor-market segmentation, and it is consistent with allocative frictions that limit firm growth or worker movement toward better matches. It is not, by itself, a direct test of misallocation. The GEIH does not identify firms or observe value added, capital, or marginal products, so the paper cannot determine whether the remaining premium reflects unobserved worker sorting or firm-side wage differences, nor whether those differences correspond to allocative inefficiencies. Further research should extend the analysis with firm-level measures of output, capital, employment and productivity to test the allocative-frictions interpretation linked to the household-survey wage patterns documented here.


\appendix
\section{Additional Tables and Figures}
\label{app:additional-descriptives}

This appendix reports the detailed employment, income, wage-growth, and
worker-characteristics tables that support the compact exhibits in the main
text. It also reports median real hourly wages by firm size and worker group for
2008 and 2025, 2025 wage gaps by worker group, plus supporting descriptive
figures and regression robustness exercises.

\subsection{Detailed Descriptive Tables}

Table~\ref{tab:descriptive-employment-size} reports the pooled employment
distribution by firm size. Tables~\ref{tab:descriptive-income-size-comparison}
and \ref{tab:descriptive-worker-characteristics-comparison} compare the pooled
sample with 2025. Tables~\ref{tab:descriptive-income-size} and
\ref{tab:descriptive-worker-characteristics} report the pooled income and worker
composition profiles by firm size in more detail. Tables
\ref{tab:descriptive-income-size-2025} and
\ref{tab:descriptive-worker-characteristics-2025} repeat the same exercise for
2025. Table~\ref{tab:descriptive-wage-growth-2008-2025} reports wage changes
between 2008 and 2025, Table~\ref{tab:descriptive-wage-groups-mean-2008-2025}
reports the corresponding mean wages by worker group, Table~\ref{tab:descriptive-wage-groups-median-2008-2025}
reports medians, and
Table~\ref{tab:descriptive-wage-gaps-2025} reports 2025 wage gaps by sex,
formality status, education group, and sector.
The appendix tables reinforce the same message as the main text. The raw
firm-size premium is visible in means, quartiles, and medians; the 2025-only
tables look similar to the pooled tables; and the composition tables show that
larger firms are simultaneously more formal and more intensive in higher
education. The wage-gap table adds that formality, education, and sector
differences are large even before conditioning on firm size, so observable
sorting is an important part of the empirical environment.

\begin{table}[p]
  \centering
  \caption{Mean real hourly labor income gaps by worker group, 2025}
  \label{tab:descriptive-wage-gaps-2025}
  \scriptsize
  \begin{tabular}{lp{0.38\textwidth}rrr}
    \toprule
    Dimension & Group & Employment share & Mean income & Gap \\
    \midrule
    Sex & Men & 54.3\% & 12,849 & 1.6\% \\
    Sex & Women & 45.7\% & 12,408 & -1.9\% \\
    Formality & Formal & 57.2\% & 16,571 & 31.0\% \\
    Formality & Informal & 42.8\% & 7,401 & -41.5\% \\
    Education & Preschool or less & 0.9\% & 4,497 & -64.4\% \\
    Education & Primary & 10.6\% & 6,282 & -50.3\% \\
    Education & Secondary & 44.3\% & 7,751 & -38.7\% \\
    Education & Higher education & 44.3\% & 19,214 & 51.9\% \\
    Sector & T - Household employers & 2.9\% & 7,183 & -43.2\% \\
    Sector & I - Accommodation and food & 7.8\% & 7,737 & -38.8\% \\
    Sector & H - Transport & 8.8\% & 8,524 & -32.6\% \\
    Sector & F - Construction & 6.9\% & 9,385 & -25.8\% \\
    Sector & G - Commerce & 19.4\% & 9,808 & -22.4\% \\
    Sector & R/S - Arts and other services & 6.1\% & 9,999 & -20.9\% \\
    Sector & C - Manufacturing & 12.8\% & 11,010 & -12.9\% \\
    Sector & D/E - Utilities and water & 1.6\% & 12,795 & 1.2\% \\
    Sector & A - Agriculture & 0.9\% & 13,975 & 10.5\% \\
    Sector & L/M/N - Business services & 12.2\% & 15,104 & 19.4\% \\
    Sector & Q - Health & 5.6\% & 16,668 & 31.8\% \\
    Sector & P - Education & 5.0\% & 20,772 & 64.2\% \\
    Sector & K - Finance and insurance & 2.7\% & 22,241 & 75.9\% \\
    Sector & J - Information and communication & 2.6\% & 23,057 & 82.3\% \\
    Sector & O - Public administration & 4.4\% & 23,893 & 88.9\% \\
    Sector & B - Mining & 0.3\% & 30,685 & 142.6\% \\
    \bottomrule
  \end{tabular}
  \vspace{0.3em}
  \begin{minipage}{0.95\textwidth}
  \footnotesize
  Notes: Statistics use 2025 observations and GEIH expansion weights. Employment shares are computed within each dimension. Mean income is real hourly labor income in constant 2025 pesos. The gap is the percent difference in each group's weighted mean relative to the overall 2025 weighted mean. Formality rows exclude pensioned occupied workers. Sector rows omit extraterritorial organizations.
  \end{minipage}
\end{table}



\begin{table}[htbp]
  \centering
  \caption{Employment distribution by firm size}
  \label{tab:descriptive-employment-size}
  \small
  \begin{tabular}{lrrrrr}
    \toprule
    Firm size & Observations & Worker-years (m) & Pooled share & 2008 share & 2025 share \\
    \midrule
    Solo & 1,435,259 & 50.6 & 31.1\% & 34.5\% & 28.9\% \\
    2-3 & 494,122 & 19.6 & 12.0\% & 16.6\% & 9.8\% \\
    4-5 & 236,560 & 9.9 & 6.1\% & 6.8\% & 6.2\% \\
    6-10 & 213,420 & 10.1 & 6.2\% & 5.9\% & 6.3\% \\
    11-19 & 141,540 & 6.8 & 4.2\% & 3.7\% & 4.6\% \\
    20-30 & 128,837 & 6.5 & 4.0\% & 3.7\% & 4.1\% \\
    31-50 & 105,125 & 5.5 & 3.4\% & 3.0\% & 3.8\% \\
    51-100 & 111,905 & 6.1 & 3.8\% & 3.5\% & 4.4\% \\
    101+ & 985,554 & 47.5 & 29.2\% & 22.2\% & 31.9\% \\
    \bottomrule
  \end{tabular}
  \vspace{0.3em}
  \begin{minipage}{0.95\textwidth}
  \footnotesize
  Notes: Worker-years and employment shares use GEIH expansion weights. The sample matches the baseline regression sample.
  \end{minipage}
\end{table}


\begin{table}[htbp]
  \centering
  \caption{Real hourly labor income by firm size: pooled sample and 2025}
  \label{tab:descriptive-income-size-comparison}
  \small
  \begin{tabular}{lrrrr}
    \toprule
    Firm size & Pooled mean & 2025 mean & Pooled premium & 2025 premium \\
    \midrule
    Solo & 7,248 & 8,351 & 0.0\% & 0.0\% \\
    2-3 & 7,959 & 8,830 & 9.8\% & 5.7\% \\
    4-5 & 8,912 & 9,503 & 23.0\% & 13.8\% \\
    6-10 & 9,608 & 10,609 & 32.6\% & 27.0\% \\
    11-19 & 10,852 & 12,346 & 49.7\% & 47.8\% \\
    20-30 & 10,782 & 11,997 & 48.8\% & 43.6\% \\
    31-50 & 11,576 & 13,428 & 59.7\% & 60.8\% \\
    51-100 & 12,697 & 14,481 & 75.2\% & 73.4\% \\
    101+ & 15,735 & 18,502 & 117.1\% & 121.5\% \\
    \bottomrule
  \end{tabular}
  \vspace{0.3em}
  \begin{minipage}{0.95\textwidth}
  \footnotesize
  Notes: Pooled statistics use 2008--2019 and 2021--2025. All income values are real hourly labor income in constant 2025 pesos and use GEIH expansion weights. Premiums are percent differences in weighted means relative to solo workers within each column.
  \end{minipage}
\end{table}


\begin{table}[htbp]
  \centering
  \caption{Worker composition by firm size: pooled sample and 2025}
  \label{tab:descriptive-worker-characteristics-comparison}
  \small
  \begin{tabular}{lrrrrrr}
    \toprule
    & \multicolumn{2}{c}{Women} & \multicolumn{2}{c}{Formal} & \multicolumn{2}{c}{Higher education} \\
    \cmidrule(lr){2-3} \cmidrule(lr){4-5} \cmidrule(lr){6-7}
    Firm size & Pooled & 2025 & Pooled & 2025 & Pooled & 2025 \\
    \midrule
    Solo & 43.6\% & 41.3\% & 6.7\% & 8.3\% & 13.1\% & 16.7\% \\
    2-3 & 33.2\% & 34.6\% & 11.3\% & 13.0\% & 14.2\% & 17.0\% \\
    4-5 & 32.5\% & 36.5\% & 22.1\% & 24.2\% & 18.3\% & 20.2\% \\
    6-10 & 36.4\% & 40.5\% & 41.6\% & 46.6\% & 26.1\% & 29.9\% \\
    11-19 & 39.2\% & 40.4\% & 60.7\% & 64.4\% & 35.0\% & 38.8\% \\
    20-30 & 40.6\% & 42.7\% & 74.6\% & 81.2\% & 40.4\% & 45.4\% \\
    31-50 & 40.3\% & 41.2\% & 84.2\% & 90.4\% & 44.2\% & 51.0\% \\
    51-100 & 40.8\% & 43.7\% & 89.5\% & 94.3\% & 47.1\% & 53.4\% \\
    101+ & 44.8\% & 45.8\% & 94.8\% & 98.0\% & 58.4\% & 62.9\% \\
    \bottomrule
  \end{tabular}
  \vspace{0.3em}
  \begin{minipage}{0.95\textwidth}
  \footnotesize
  Notes: Pooled statistics use 2008--2019 and 2021--2025. All statistics use GEIH expansion weights. Higher education corresponds to workers classified as superior or university educated.
  \end{minipage}
\end{table}


\begin{table}[htbp]
  \centering
  \caption{Real hourly labor income by firm size, pooled years}
  \label{tab:descriptive-income-size}
  \small
  \begin{tabular}{lrrrrrr}
    \toprule
    Firm size & Mean & P25 & Median & P75 & S.D. & Raw premium \\
    \midrule
    Solo & 5,690 & 2,366 & 3,943 & 6,169 & 11,205 & 0.0\% \\
    2-3 & 6,464 & 3,067 & 4,784 & 6,779 & 12,543 & 13.6\% \\
    4-5 & 7,444 & 3,943 & 5,430 & 7,163 & 21,476 & 30.8\% \\
    6-10 & 8,403 & 4,647 & 5,922 & 7,854 & 16,634 & 47.7\% \\
    11-19 & 9,746 & 5,204 & 6,383 & 8,774 & 17,881 & 71.3\% \\
    20-30 & 9,972 & 5,468 & 6,669 & 9,341 & 21,069 & 75.3\% \\
    31-50 & 10,744 & 5,714 & 6,939 & 10,321 & 16,435 & 88.8\% \\
    51-100 & 11,612 & 5,809 & 7,176 & 11,255 & 18,818 & 104.1\% \\
    101+ & 14,650 & 6,166 & 8,624 & 16,169 & 20,771 & 157.5\% \\
    \bottomrule
  \end{tabular}
  \vspace{0.3em}
  \begin{minipage}{0.95\textwidth}
  \footnotesize
  Notes: Incomes are pooled over 2008--2019 and 2021--2025 and expressed as real hourly labor income in constant 2025 pesos. All statistics are weighted by GEIH expansion weights. The raw premium is the percent difference in the weighted mean relative to solo workers.
  \end{minipage}
\end{table}


\begin{table}[htbp]
  \centering
  \caption{Worker characteristics by firm size, pooled years}
  \label{tab:descriptive-worker-characteristics}
  \small
  \begin{tabular}{lrrrr}
    \toprule
    Firm size & Women & Formal & Higher education & Mean log wage \\
    \midrule
    Solo & 49.7\% & 11.1\% & 19.9\% & 8.457 \\
    2-3 & 41.5\% & 17.1\% & 22.9\% & 8.596 \\
    4-5 & 40.2\% & 30.9\% & 26.7\% & 8.744 \\
    6-10 & 42.1\% & 50.0\% & 33.5\% & 8.847 \\
    11-19 & 43.8\% & 67.0\% & 41.6\% & 8.949 \\
    20-30 & 44.0\% & 78.8\% & 45.8\% & 8.986 \\
    31-50 & 43.2\% & 86.8\% & 48.9\% & 9.049 \\
    51-100 & 43.7\% & 91.2\% & 52.3\% & 9.114 \\
    101+ & 45.9\% & 95.6\% & 62.2\% & 9.300 \\
    \bottomrule
  \end{tabular}
  \vspace{0.3em}
  \begin{minipage}{0.95\textwidth}
  \footnotesize
  Notes: Statistics are pooled over 2008--2019 and 2021--2025 and weighted by GEIH expansion weights. Higher education corresponds to workers classified as superior or university educated.
  \end{minipage}
\end{table}


\begin{table}[htbp]
  \centering
  \caption{Real hourly labor income by firm size, 2025}
  \label{tab:descriptive-income-size-2025}
  \small
  \begin{tabular}{lrrrrrr}
    \toprule
    Firm size & Mean & P25 & Median & P75 & S.D. & Raw premium \\
    \midrule
    Solo & 8,351 & 3,606 & 5,769 & 8,942 & 11,807 & 0.0\% \\
    2-3 & 8,830 & 4,615 & 6,522 & 8,654 & 13,608 & 5.7\% \\
    4-5 & 9,503 & 5,440 & 6,923 & 8,654 & 16,211 & 13.8\% \\
    6-10 & 10,609 & 6,393 & 7,466 & 9,532 & 17,310 & 27.0\% \\
    11-19 & 12,346 & 6,844 & 7,867 & 11,276 & 18,438 & 47.8\% \\
    20-30 & 11,997 & 7,139 & 8,145 & 11,538 & 13,414 & 43.6\% \\
    31-50 & 13,428 & 7,141 & 8,515 & 13,516 & 17,573 & 60.8\% \\
    51-100 & 14,481 & 7,212 & 8,897 & 14,685 & 18,799 & 73.4\% \\
    101+ & 18,502 & 7,615 & 11,037 & 20,769 & 21,986 & 121.5\% \\
    \bottomrule
  \end{tabular}
  \vspace{0.3em}
  \begin{minipage}{0.95\textwidth}
  \footnotesize
  Notes: Incomes use only 2025 observations and are expressed as real hourly labor income in constant 2025 pesos. All statistics are weighted by GEIH expansion weights. The raw premium is the percent difference in the weighted mean relative to solo workers in 2025.
  \end{minipage}
\end{table}


\begin{table}[htbp]
  \centering
  \caption{Worker characteristics by firm size, 2025}
  \label{tab:descriptive-worker-characteristics-2025}
  \small
  \begin{tabular}{lrrrr}
    \toprule
    Firm size & Women & Formal & Higher education & Mean log wage \\
    \midrule
    Solo & 46.2\% & 14.1\% & 25.2\% & 8.639 \\
    2-3 & 42.4\% & 20.1\% & 26.1\% & 8.769 \\
    4-5 & 42.6\% & 32.2\% & 28.3\% & 8.879 \\
    6-10 & 44.0\% & 54.4\% & 36.3\% & 9.009 \\
    11-19 & 44.2\% & 70.4\% & 46.1\% & 9.128 \\
    20-30 & 45.1\% & 84.2\% & 49.5\% & 9.168 \\
    31-50 & 44.2\% & 91.9\% & 56.3\% & 9.254 \\
    51-100 & 46.8\% & 94.8\% & 58.4\% & 9.311 \\
    101+ & 47.5\% & 98.3\% & 67.6\% & 9.511 \\
    \bottomrule
  \end{tabular}
  \vspace{0.3em}
  \begin{minipage}{0.95\textwidth}
  \footnotesize
  Notes: Statistics use only 2025 observations and are weighted by GEIH expansion weights. Higher education corresponds to workers classified as superior or university educated.
  \end{minipage}
\end{table}


\begin{table}[!htbp]
  \centering
  \caption{Growth in mean real hourly labor income by firm size, 2008--2025}
  \label{tab:descriptive-wage-growth-2008-2025}
  \small
  \begin{tabular}{lrrr}
    \toprule
    Firm size & 2008 mean & 2025 mean & Change \\
    \midrule
    All firm sizes & 8,736 & 12,647 & 44.8\% \\
    Solo & 6,100 & 8,351 & 36.9\% \\
    2-3 & 7,223 & 8,830 & 22.2\% \\
    4-5 & 8,291 & 9,503 & 14.6\% \\
    6-10 & 8,249 & 10,609 & 28.6\% \\
    11-19 & 9,193 & 12,346 & 34.3\% \\
    20-30 & 9,347 & 11,997 & 28.3\% \\
    31-50 & 10,960 & 13,428 & 22.5\% \\
    51-100 & 11,193 & 14,481 & 29.4\% \\
    101+ & 13,369 & 18,502 & 38.4\% \\
    \bottomrule
  \end{tabular}
  \vspace{0.3em}
  \begin{minipage}{0.95\textwidth}
  \footnotesize
  Notes: Values are weighted mean real hourly labor income in constant 2025 pesos. Change is the percent change in the weighted mean between 2008 and 2025.
  \end{minipage}
\end{table}


\begin{table}[!htbp]
  \centering
  \caption{Mean real hourly labor income by firm size and worker group, 2008 and 2025}
  \label{tab:descriptive-wage-groups-mean-2008-2025}
  \small
  \begin{tabular}{lrrrrr}
    \toprule
    Firm size & All & Men & Women & Formal & Informal \\
    \midrule
    \multicolumn{6}{l}{\textit{Panel A: 2008}} \\
    All firm sizes & 8,736 & 9,087 & 8,264 & 12,476 & 6,367 \\
    Solo & 6,100 & 6,790 & 5,357 & 12,658 & 5,626 \\
    2-3 & 7,223 & 7,355 & 6,980 & 11,377 & 6,674 \\
    4-5 & 8,291 & 8,680 & 7,585 & 11,831 & 7,204 \\
    6-10 & 8,249 & 8,788 & 7,399 & 10,061 & 6,890 \\
    11-19 & 9,193 & 9,807 & 8,372 & 10,187 & 7,796 \\
    20-30 & 9,347 & 9,406 & 9,259 & 10,256 & 7,261 \\
    31-50 & 10,960 & 10,654 & 11,401 & 10,719 & 11,901 \\
    51-100 & 11,193 & 10,589 & 12,111 & 11,911 & 7,075 \\
    101+ & 13,369 & 13,447 & 13,270 & 13,728 & 9,622 \\
    \addlinespace
    \multicolumn{6}{l}{\textit{Panel B: 2025}} \\
    All firm sizes & 12,647 & 12,849 & 12,408 & 16,571 & 7,401 \\
    Solo & 8,351 & 8,749 & 7,889 & 17,217 & 6,894 \\
    2-3 & 8,830 & 8,828 & 8,832 & 14,077 & 7,511 \\
    4-5 & 9,503 & 9,561 & 9,424 & 13,692 & 7,517 \\
    6-10 & 10,609 & 10,824 & 10,335 & 12,844 & 7,947 \\
    11-19 & 12,346 & 12,833 & 11,733 & 13,656 & 9,227 \\
    20-30 & 11,997 & 11,957 & 12,045 & 12,521 & 9,206 \\
    31-50 & 13,428 & 13,484 & 13,358 & 13,879 & 8,337 \\
    51-100 & 14,481 & 14,992 & 13,900 & 14,598 & 12,327 \\
    101+ & 18,502 & 18,869 & 18,096 & 18,548 & 15,843 \\
    \bottomrule
  \end{tabular}
  \vspace{0.3em}
  \begin{minipage}{0.95\textwidth}
  \footnotesize
  Notes: Values are weighted mean real hourly labor income in constant 2025 pesos. The first row pools all firm-size categories. Men and women are defined from the worker's reported sex; formal and informal are defined using the baseline formality indicator.
  \end{minipage}
\end{table}


\begin{table}[!htbp]
  \centering
  \caption{Median real hourly labor income by firm size and worker group, 2008 and 2025}
  \label{tab:descriptive-wage-groups-median-2008-2025}
  \small
  \begin{tabular}{lrrrrr}
    \toprule
    Firm size & All & Men & Women & Formal & Informal \\
    \midrule
    \multicolumn{6}{l}{\textit{Panel A: 2008}} \\
    All firm sizes & 5,128 & 5,244 & 5,034 & 7,132 & 4,195 \\
    Solo & 4,027 & 4,195 & 3,596 & 6,616 & 3,776 \\
    2-3 & 4,489 & 4,720 & 4,195 & 5,873 & 4,195 \\
    4-5 & 4,840 & 4,894 & 4,840 & 5,422 & 4,661 \\
    6-10 & 5,244 & 5,244 & 5,244 & 5,802 & 4,840 \\
    11-19 & 5,559 & 5,454 & 5,664 & 6,162 & 5,003 \\
    20-30 & 5,929 & 5,768 & 6,293 & 6,413 & 5,034 \\
    31-50 & 6,075 & 5,873 & 6,293 & 6,293 & 4,952 \\
    51-100 & 6,293 & 6,293 & 6,678 & 6,670 & 5,034 \\
    101+ & 8,055 & 7,866 & 8,390 & 8,390 & 5,387 \\
    \addlinespace
    \multicolumn{6}{l}{\textit{Panel B: 2025}} \\
    All firm sizes & 7,805 & 7,805 & 7,821 & 9,615 & 5,769 \\
    Solo & 5,769 & 6,000 & 5,538 & 10,256 & 5,385 \\
    2-3 & 6,522 & 6,593 & 6,410 & 8,395 & 5,769 \\
    4-5 & 6,923 & 6,923 & 6,923 & 8,145 & 6,250 \\
    6-10 & 7,466 & 7,466 & 7,466 & 8,077 & 6,731 \\
    11-19 & 7,867 & 7,867 & 7,867 & 8,308 & 6,923 \\
    20-30 & 8,145 & 8,027 & 8,212 & 8,308 & 7,023 \\
    31-50 & 8,515 & 8,392 & 8,654 & 8,654 & 7,133 \\
    51-100 & 8,897 & 8,654 & 9,135 & 9,030 & 7,692 \\
    101+ & 11,037 & 10,714 & 11,538 & 11,037 & 7,692 \\
    \bottomrule
  \end{tabular}
  \vspace{0.3em}
  \begin{minipage}{0.95\textwidth}
  \footnotesize
  Notes: Values are weighted median real hourly labor income in constant 2025 pesos. The first row pools all firm-size categories. Men and women are defined from the worker's reported sex; formal and informal are defined using the baseline formality indicator.
  \end{minipage}
\end{table}



\subsection{Additional Descriptive Figures}

Figures~\ref{fig:descriptive-employment-share-time},
\ref{fig:descriptive-raw-premium-time}, and
\ref{fig:descriptive-worker-composition} report the full time path of
employment shares, raw firm-size wage premia, and worker composition.
The time-series figures show persistence rather than a one-year anomaly:
employment remains concentrated at the bottom, raw wage premia remain positive,
and formality and higher education remain strongly increasing in firm size.

\begin{figure}[htbp]
  \centering
  \includegraphics[width=0.9\textwidth]{fig61_\figurelang.png}
  \caption{Employment share by firm size over time}
  \label{fig:descriptive-employment-share-time}
\end{figure}

\begin{figure}[htbp]
  \centering
  \includegraphics[width=0.9\textwidth]{fig62_\figurelang.png}
  \caption{Raw firm-size wage premium over time, relative to solo workers}
  \label{fig:descriptive-raw-premium-time}
\end{figure}

\begin{figure}[htbp]
  \centering
  \includegraphics[width=0.9\textwidth]{fig63_\figurelang.png}
  \caption{Worker composition by firm size}
  \label{fig:descriptive-worker-composition}
\end{figure}

\clearpage

\subsection{Additional Regression Tables and Figures}
\label{app:additional-regression-figures}

Figure~\ref{fig:firm-size-beta-coefficients} reports the log-point coefficients
from the full specification, and Figure~\ref{fig:firm-size-premium-trend-test}
visualizes the corresponding trend-test estimates.
Together, these figures show that the conditional firm-size coefficients are
positive in the full specification, while the estimated time trends are positive
but too imprecise to support a strong claim that the premium has systematically
steepened.

Table~\ref{tab:firm-size-local-geography-robustness} reports a local-market
robustness exercise. Because the harmonized file retains department/Bogot\'a
rather than a consistent city identifier, the table uses department as the local
geography. The most demanding column absorbs sector-department-year fixed
effects, so the firm-size coefficients are identified within the same sector,
department, and year. The large-firm coefficient remains positive in this
specification: workers in firms with 101 or more employees earn about 31 percent
more than comparable solo workers within the same sector-department-year cell.

\begin{table}[htbp]
  \centering
  \caption{Firm-size wage premium robustness with local sector-year fixed effects}
  \label{tab:firm-size-local-geography-robustness}
  \small
  \begin{tabular}{lccc}
    \toprule
    & \multicolumn{3}{c}{Dependent variable: log real hourly labor income} \\
    \cmidrule(lr){2-4}
    & (1) & (2) & (3) \\
    \midrule
    Firm size: 2-3 & 0.145*** & 0.126*** & 0.127*** \\
     & (0.014) & (0.014) & (0.015) \\
    Firm size: 4-5 & 0.234*** & 0.216*** & 0.215*** \\
     & (0.020) & (0.020) & (0.021) \\
    Firm size: 6-10 & 0.234*** & 0.215*** & 0.215*** \\
     & (0.023) & (0.023) & (0.024) \\
    Firm size: 11-19 & 0.224*** & 0.209*** & 0.209*** \\
     & (0.022) & (0.022) & (0.023) \\
    Firm size: 20-30 & 0.190*** & 0.177*** & 0.177*** \\
     & (0.022) & (0.022) & (0.023) \\
    Firm size: 31-50 & 0.195*** & 0.184*** & 0.183*** \\
     & (0.022) & (0.023) & (0.023) \\
    Firm size: 51-100 & 0.215*** & 0.203*** & 0.203*** \\
     & (0.022) & (0.022) & (0.023) \\
    Firm size: 101+ & 0.282*** & 0.272*** & 0.268*** \\
     & (0.030) & (0.032) & (0.032) \\
    \midrule
    Worker and formality controls & Yes & Yes & Yes \\
    Fixed effects & $s\times t$ & $s\times t$, $d\times t$ & $s\times d\times t$ \\
    Observations & 3,852,322 & 3,852,322 & 3,852,211 \\
    $R^2$ & 0.362 & 0.374 & 0.380 \\
    \bottomrule
  \end{tabular}
  \vspace{0.3em}
  \begin{minipage}{0.95\textwidth}
  \footnotesize
  Notes: The omitted category is solo workers. The sample is restricted to observations with nonmissing department codes. Worker controls include a female-worker indicator, age, age squared, education dummies, and labor formality. The local geography $d$ is department/Bogot\'a, because the harmonized file does not retain a consistent city or municipality identifier. Column (3) is the most stringent specification and absorbs sector-department-year fixed effects, so identification comes from wage differences across firm-size categories within the same sector, department, and year. Standard errors are clustered by sector-department cells. Significance levels: * $p<0.10$, ** $p<0.05$, *** $p<0.01$.
  \end{minipage}
\end{table}


\begin{figure}[htbp]
  \centering
  \includegraphics[width=0.9\textwidth]{fig58_\figurelang.png}
  \caption{Firm-size coefficients in the full specification}
  \label{fig:firm-size-beta-coefficients}
\end{figure}

\begin{figure}[htbp]
  \centering
  \includegraphics[width=0.9\textwidth]{fig60_\figurelang.png}
  \caption{Trend test for the firm-size wage premium}
  \label{fig:firm-size-premium-trend-test}
\end{figure}


\bibliographystyle{apalike}
\bibliography{Paper/references}

\end{document}
